\documentclass[11pt]{article}
\usepackage[margin=1.15in]{geometry}
\usepackage{amsmath,amssymb,amsthm}
\usepackage{mathpazo}
\usepackage[T1]{fontenc}
\usepackage{microtype}
\usepackage{setspace}
\usepackage{enumitem}
\usepackage{array}
\usepackage[colorlinks=true,linkcolor=black,citecolor=black,urlcolor=black]{hyperref}
\onehalfspacing
\setlength{\parindent}{1.5em}
\setlength{\parskip}{0.2em}

\newcommand{\F}{\mathcal F}
\newcommand{\Q}{\mathcal Q}
\newcommand{\C}{\mathcal C}
\newcommand{\PP}{\mathbb P}
\newcommand{\E}{\mathbb E}
\newcommand{\feas}{\mathrm{feas}}
\newcommand{\adm}{\mathrm{adm}}
\newcommand{\app}{\mathrm{app}}
\newcommand{\WB}{\mathsf{WB}}

\theoremstyle{plain}
\newtheorem{proposition}{Proposition}
\theoremstyle{definition}
\newtheorem{definition}{Definition}
\newtheorem{remark}{Remark}

\title{Deliberate Ignorance and the Law:\\Where Philosophical and Doctrinal Willful Blindness Diverge}
\author{Flyxion\\Independent Researcher}
\date{October 2026}

\begin{document}
\maketitle

\begin{abstract}
\noindent Willful blindness in the criminal law and deliberate ignorance in a philosophical account of responsibility are not the same concept, and treating them as interchangeable obscures what each can do. This paper sets the doctrine as stated by the United States Supreme Court in \emph{Global-Tech Appliances v. SEB} and by the Supreme Court of Canada in \emph{Sansregret}, \emph{Jorgensen}, and \emph{Briscoe} against a framework in which responsibility attaches to the admissible acquisition of information. It distinguishes three layers that discussions of the subject tend to run together: philosophical responsibility for the establishment of an informational domain, the legal knowledge requirements that substitute willful blindness for knowledge, and the evidentiary standards by which a deliberate avoidance is proved. It shows that the subjective suspicion required by the doctrine and the evidential indication on which the philosophical account relies are logically independent, so that neither implies the other, and that the law's exclusion of negligence corresponds to a division of labor between doctrines that the philosophical account unifies. It formalizes the difference among recklessness, negligence, and willful blindness through the two stages of a decision, analyzes the partial or sham inquiry that Canadian and American cases treat differently, and considers organizations, where the philosophical notion of institutional ignorance without individual blindness meets statutory rules of attribution. Authorities are marked according to whether their text was read directly or only through secondary sources, and the paper offers no legal advice.
\end{abstract}

\section{Introduction}

The phrase willful blindness is used in at least three ways, and the differences among them are easy to lose. It names a philosophical failure, that of an agent who arranges not to learn what the agent suspects would be unwelcome. It names a legal device, a rule that substitutes deliberate avoidance of knowledge for knowledge where knowledge is an element of an offence or a civil wrong. And it names a problem of proof, that of showing from the circumstances of a case that an accused or defendant did in fact avoid learning something, since a state of mind is not directly observable. Each of these has its own standards, its own consequences, and its own reasons for being what it is. A discussion that moves among them without notice will find that a point proved in one register does not carry over to another, and that apparent disagreements between lawyers and philosophers are sometimes disagreements about different things.

The parent essay in this series, \emph{The Domain and Its Establishment} \cite{flyxion2026domain}, develops an account in which an agent is answerable not only for operating within a domain of information but also for how that domain was established, maintained, or deliberately restricted. Its central distinction is between the information an agent possesses and the information the agent could reasonably have acquired, and within the second it distinguishes deficiencies that are merely negligent from those that involve a motivational component of avoidance. The first companion paper, \emph{Correct Relative to Its Inputs} \cite{flyxion2026inputs}, applies the account to automated decisions and to organizations. The present paper turns to the law, where the concept of deliberate ignorance has a long doctrinal history, and asks how the framework and the doctrine relate. The question is not how the law should be changed, and it is not how the framework should be revised to match the law. It is where the two coincide, where they diverge, and why.

The claim defended is that the law's willful blindness occupies a narrow region of the framework and also extends beyond that region in one direction, so that neither contains the other. The doctrine is narrow because it excludes negligence by design: the courts that apply it are explicit that a defendant who should have known is not for that reason a defendant who is blind. It extends beyond the framework's account of deliberate avoidance because it turns on what the defendant actually suspected, which need not match what the evidence available to the defendant indicated. The division between possessed and acquirable information that is central to the framework reappears in the law as a division between doctrines: willful blindness addresses the defendant's own state of mind and deliberate avoidance, while the law of negligence, due diligence, and regulatory duty addresses what the defendant should have found out. The framework treats the two as parts of one account of responsibility for the establishment of a domain, whereas the law keeps them apart because their consequences differ.

The paper proceeds as follows. Section~2 recalls the framework and adds the two notions that the comparison needs, the credence of an agent and the avoidance act. Sections~3 and~4 state the doctrine in the United States and in Canada. Section~5 separates the three layers of philosophical responsibility, legal knowledge, and evidential proof. Section~6 establishes the independence of subjective suspicion and evidential indication, and Section~7 collects the further divergences. Section~8 recasts recklessness, negligence, and willful blindness as positions in the two-stage structure of a decision, Section~9 analyzes partial and sham inquiries, and Section~10 takes up evidence and hindsight. Section~11 considers organizations, Section~12 carries one hypothetical through five variants, and Section~13 asks whether the law ought to reach further. Section~14 states limits and conclusions, and an appendix records which authorities were read in the original and which were consulted only through secondary accounts.

\paragraph{Scope and caution.} Only federal law of the United States and the criminal and public law of Canada are considered, and within them only the doctrines that bear on deliberate ignorance. The paper is a work of legal and philosophical analysis and is not legal advice. The statements of doctrine are drawn from the sources identified in Appendix~A and are limited to what those sources support, and the reader who needs the law applied to a particular case should consult the primary authorities and current statutory text.

\section{The framework and two additions}

The framework is recalled here in compressed form so that the paper can be read independently, and its development is in \cite{flyxion2026domain}. Let $(\Omega,\F,\PP)$ be a probability space with a filtration $(\F_t)$, where $\F^{(i)}_t$ is the information agent $i$ actually possesses at time $t$. A decision proceeds in two stages. In the first, the agent chooses an investigation $q$ from the set $\Q_t$ of acquisition procedures, using only $\F^{(i)}_t$; in the second, after the observation $Z_q$ arrives, the agent chooses an action using the enlarged state $\F^{(i)}_t\vee\sigma(Z_q)$. The feasible investigations $\Q_t^{\feas}$ are those the agent can carry out, the applicable investigations $\Q_t^{\app}$ are those the governing standard of care calls for, and the admissible investigations are $\Q_t^{\adm}=\Q_t^{\feas}\cap\Q_t^{\app}$. An omission of an admissible investigation that was reasonably indicated by the evidence is deficient, with a violation measure $\nu_t>0$, and a deficient omission is of type D2 when a motivational component is present, namely that the avoidance is attributable to a suspicion of what the investigation would show, and of type D1 otherwise. Types of omission that are not deficient include the unavoidable (U), the not indicated (N), and the justified by a competing constraint (J).

The comparison with the law requires two additions. The first is the distinction between what the evidence supports and what the agent believes. Let $f\in\F$ be a fact, an event whose truth is material to the decision or is an element of a legal wrong, and let $q_f$ be an investigation that would settle it. Define the \emph{evidence-supported probability}
\[
p^{(i)}_t(f)=\PP\bigl(f\mid\F^{(i)}_t\bigr),
\]
which is a function of the information the agent holds under the standard model, and the investigation $q_f$ is \emph{indicated} when $p^{(i)}_t(f)\ge\theta_t$, where $\theta_t$ is the diligence threshold fixed by the standard of care. Define also the agent's \emph{credence} $\beta^{(i)}_t(f)\in[0,1]$, the probability the agent in fact assigns to $f$. The credence is a function of how the agent has processed the information, and it need not equal $p^{(i)}_t(f)$: an agent may overlook evidence, overread it, or hold a suspicion that the evidence does not support. The parent essay distinguished possessed from acquirable information. The additional distinction here is between information possessed and information processed, and it turns out to be the distinction on which the legal doctrine depends.

The second addition concerns the acts by which an agent avoids learning. The parent essay treated avoidance as the omission of an admissible investigation, but the law's cases involve more than omissions. An \emph{avoidance act} with respect to $f$ is any of the following. It may be the omission of an investigation $q_f$ that would settle $f$. It may be the selection of a coarser substitute $q'$ with $\sigma(Z_{q'})\subsetneq\sigma(Z_{q_f})$, so that the information obtained is less revealing about $f$ than an available alternative of comparable cost. Or it may be interference, a step that restricts what flows into the investigation or reaches the agent, such as withholding from a consultant the fact that would make the consultant's inquiry informative, or destroying or declining to receive a record. Write $\mathrm{Av}^{(i)}_t(f)$ for the set of avoidance acts available to or performed by agent $i$ with respect to $f$.

\section{The doctrine in the United States}

The leading statement is the Supreme Court's opinion in \emph{Global-Tech Appliances, Inc.\ v.\ SEB S.A.}, 563 U.S.\ 754 (2011), written by Justice Alito for an eight-member majority. The case concerned the civil provision on induced patent infringement, 35 U.S.C.\ \S\,271(b), and the Court held that inducement requires knowledge that the induced acts constitute infringement. The question was then whether the knowledge requirement could be satisfied by willful blindness, and the Court answered that it could. Its reasoning drew on the criminal law, where the doctrine was described as well established, and it found no reason to exclude it from civil patent law (at 766 to 768). The Court traced the doctrine to \emph{United States v.\ Jewell}, 532 F.2d 697 (9th Cir.\ 1976) (en banc), noted the definition of knowledge in the Model Penal Code's proposed draft as including awareness of a high probability of a fact's existence unless the person actually believes it does not exist, and observed that every court of appeals had embraced the doctrine with the possible exception of the District of Columbia Circuit (at 766 to 769 and n.\,9).

The Court stated two requirements, and these give the American formulation its limited scope. First, the defendant must subjectively believe that there is a high probability that a fact exists. Second, the defendant must take deliberate actions to avoid learning of that fact (at 769). The Court expressly distinguished the doctrine from recklessness and from negligence. A reckless defendant, in the Court's description, is one who merely knows of a substantial and unjustified risk of wrongdoing, and a negligent defendant is one who should have known of a similar risk but in fact did not (at 770). The Federal Circuit's standard of deliberate indifference to a known risk was held to depart from the proper standard in two respects, of which the first is that it permitted a finding of knowledge where there was only a known risk (at 770). The Court's reasoning treats the high-probability belief and the deliberate avoidance together as what distinguishes a defendant who is treated as knowing from one who is merely at fault.

The facts illustrate the second requirement. The defendant copied a competitor's fryer, bought in a foreign market where the product bore no United States patent marking, and commissioned a right-to-use opinion from a patent attorney without telling the attorney that the product had been copied. The Court held the evidence more than sufficient for a jury to find willful blindness, reasoning that the copying showed the defendant regarded the design as valuable, that the defendant knew products made for overseas markets typically lack United States markings, and that the only apparent motive for withholding the fact of copying from the attorney was to manufacture a claim of plausible deniability (at 771). The avoidance in this case is an interference in the sense of Section~2: the defendant did not decline to consult an attorney, but arranged the consultation so that it could not reveal the fact at issue.

Justice Kennedy dissented. He agreed that the statute requires knowledge but argued that willful blindness is not knowledge, that a court should not enlarge a statutory prohibition by analogy to another body of law, and that the case should be remanded for the Federal Circuit to assess whether the evidence supported a finding of actual knowledge (at 772 to 775). The dissent matters for the present inquiry because it denies the equivalence on which the doctrine rests, and Section~13 returns to it.

\section{The doctrine in Canada}

The Canadian doctrine is stated in a line of Supreme Court of Canada decisions, of which the principal are \emph{R.\ v.\ Sansregret}, [1985] 1 S.C.R.\ 570, \emph{R.\ v.\ Jorgensen}, [1995] 4 S.C.R.\ 55, and \emph{R.\ v.\ Briscoe}, 2010 SCC 13. The primary text of these decisions could not be retrieved in the course of preparing this paper, and the account that follows relies on secondary sources identified in Appendix~A. It is therefore stated at a level of generality that those sources support, and the paragraph references given are the ones those sources report.

Willful blindness, spelled wilful blindness in Canadian usage, attributes knowledge to a person whose suspicion has been aroused to the point at which the person sees the need for further inquiry, and who deliberately chooses not to make the inquiry in order to avoid confirming what is suspected. In \emph{Briscoe}, a unanimous decision written by Justice Charron concerning an accused charged as a party to kidnapping, aggravated sexual assault, and murder, the Court treated willful blindness as a substitute for actual knowledge and not as a separate form of fault, so that where it is established the imputed knowledge is equivalent to actual subjective knowledge (reported at para.\ 21). The Court distinguished it from recklessness and from negligence, and the decision is reported to give as a reason for the limited scope of the doctrine that a wider definition would make it indistinguishable from the civil doctrine of negligence (reported at para.\ 23). \emph{Sansregret} is reported to have drawn the contrast with recklessness in the following terms: recklessness involves awareness of a risk and proceeding in the face of it, whereas willful blindness involves awareness of the need for inquiry and a deliberate decision to decline it because the accused does not want to know the truth (reported at para.\ 22). \emph{Jorgensen} is reported to have described the core of the test as a suspicion of the fact, a realization of its probability, and a refraining from confirmation in order to preserve the ability to deny knowledge, and to have emphasized that the determination is contextual (reported at pp.\ 158 to 159 and paras.\ 101 and 103).

Two further features of the Canadian law are reported in the decisions of the provincial courts of appeal. The first is that the inquiry is subjective: the question is not whether the accused should have been suspicious but whether the accused was in fact suspicious, and willful blindness is not constructive knowledge on a reasonableness standard (reported from \emph{R.\ v.\ Malfara} and \emph{R.\ v.\ Callejas}). The second concerns an inquiry that is made but is inadequate. In \emph{R.\ v.\ Lagace} the Ontario Court of Appeal is reported to have required a real suspicion that causes the accused to see the need for inquiry, and to have held that where some inquiry was made the Crown must prove beyond a reasonable doubt that the accused remained suspicious and avoided further inquiry, while an accused who took all reasonable steps to learn the truth is not willfully blind (reported at paras.\ 26 and following). The mere failure to inquire is not enough, and the ignorance must be deliberate (reported from \emph{R.\ v.\ Farmer}). Willful blindness can be proved by circumstantial evidence, and a jury instruction on it has an air of reality where the circumstances were inherently suspicious (reported from \emph{R.\ v.\ Anderson}).

In outline, therefore, the Canadian formulation shares with the American the exclusion of negligence, the treatment of the doctrine as a substitute for knowledge, and the requirement of a subjective state. It differs in its centre of gravity. The American statement speaks of a belief in a high probability and of deliberate actions to avoid learning, while the Canadian statement speaks of a suspicion that has arisen to the point of seeing a need for inquiry and of a deliberate decision not to inquire. Whether these amount to the same threshold in practice is a matter on which the sources consulted do not permit a firm conclusion, and the paper does not assume that they do.

\section{Three layers}

Discussion of deliberate ignorance goes wrong most often by treating a claim made at one layer as if it were a claim at another. The comparison is clearer if the layers are stated as functions with different arguments.

The first layer is that of philosophical responsibility for the establishment of a domain. Its output is a classification of an omission as unavoidable, not indicated, justified, or deficient, and of a deficient omission as negligent (D1) or motivated (D2). Its arguments are the information held $\F^{(i)}_t$, the feasible and applicable investigations, the standard of care with its threshold $\theta_t$, and the motivational component. It is a classification of the agent's conduct in relation to a standard, and it is made from the standpoint of an evaluator who is entitled to ask what the agent could and should have done.

The second layer is that of legal knowledge requirements. Its output is a truth value for the proposition that the element of knowledge is satisfied, and its arguments are narrower. Two formulations are relevant. Write $\WB^{\mathrm{US}}_t(i,f)$ for the proposition that agent $i$ at time $t$ satisfies the American formulation with respect to fact $f$, namely that $\beta^{(i)}_t(f)\ge\eta$ for a level $\eta$ that the law describes as a high probability, and that the agent has performed an avoidance act with respect to $f$ with the aim of not learning it. Write $\WB^{\mathrm{CA}}_t(i,f)$ for the Canadian formulation, namely that the agent's suspicion about $f$ has risen to a level $\varsigma$ at which the agent sees the need for inquiry, and that the agent has deliberately declined to inquire, or to inquire further, in order to avoid confirming it. Both formulations contain the same implicit motivational predicate $\mathrm{Mot}^{(i)}_t(f)$, that the avoidance is attributable to a suspicion of what learning would show. The levels $\eta$ and $\varsigma$ are parameters of the legal test, and they are not the diligence threshold $\theta_t$ of the standard of care.

The third layer is evidential. Its output is a finding by a trier of fact that the second-layer proposition holds, and its arguments are the evidence before the trier, the rules of inference and admissibility, and the burden and standard of proof. The existence of willful blindness in the sense of the second layer and the finding of it in the sense of the third come apart whenever the evidence is insufficient or misleading, and the possibility of error in both directions is built into the structure of adjudication and is not an anomaly.

The three layers answer different questions. The first asks whether the agent answered adequately for how the agent's informational domain came to be what it was. The second asks whether the agent's state of mind and conduct satisfy a particular legal element, so that a consequence attaches that would otherwise require proof of knowledge. The third asks whether a tribunal can be satisfied that the second is the case. A statement that is true of one layer, for instance that negligence is blameworthy, does not transfer to another, for instance to the claim that negligence supplies the knowledge element, and the law's insistence on the latter is not a denial of the former.

\section{Subjective suspicion and evidential indication}

The central structural difference between the framework and the doctrine concerns what triggers the duty or the culpability. In the framework the trigger is evidential: an investigation is indicated when the evidence in the agent's information supports the probability of the fact at or above the diligence threshold, $p^{(i)}_t(f)\ge\theta_t$. In the doctrine the trigger is subjective: the agent's own credence or suspicion must reach the level the test requires. The two are logically independent.

\begin{proposition}\label{prop:independence}
(i) Suppose the agent is a Bayesian reasoner under the standard model, so that $\beta^{(i)}_t(f)=p^{(i)}_t(f)$. Then the first element of the American formulation, $\beta^{(i)}_t(f)\ge\eta$, implies that the investigation is indicated whenever $\eta\ge\theta_t$, and the converse fails whenever $\theta_t<\eta$. (ii) For agents who are not Bayesian reasoners, each of the four combinations of indicated or not indicated with credence at or above $\eta$ or below it is realizable, and in each combination an avoidance act can be present or absent.
\end{proposition}

\begin{proof}
(i) With $\beta=p$, the condition $\beta\ge\eta$ is the condition $p\ge\eta$. If $\eta\ge\theta_t$ then $p\ge\eta\ge\theta_t$, so the investigation is indicated. If $\theta_t<\eta$, an agent with $p\in[\theta_t,\eta)$ has the investigation indicated and does not satisfy the first element, so indication does not imply the element. (ii) The credence $\beta^{(i)}_t(f)$ is a function of the agent's processing of $\F^{(i)}_t$ and is not constrained to equal $p^{(i)}_t(f)$. Fix an information state with $p\ge\theta_t$ and let the agent overlook the evidence so that $\beta<\eta$: the investigation is indicated and the credence is low. Fix an information state with $p<\theta_t$ and let the agent overread it so that $\beta\ge\eta$: the investigation is not indicated and the credence is high. The remaining two combinations are obtained by letting credence follow evidence. The presence of an avoidance act is independent of either, since it is a choice made after the credence has formed.
\end{proof}

The four combinations are worth naming because each corresponds to a recognizable case. An agent with indicated evidence and a high credence who avoids learning is the paradigm of the overlap, deliberate avoidance in both senses, and it is the case that both the framework and the doctrine call willful blindness. An agent with indicated evidence and a low credence who does not inquire is negligent in the framework's sense (type D1) and is outside the doctrine, since the defendant did not in fact believe the fact probable. An agent with no indicating evidence and a high credence who avoids learning is outside the framework's duty, since the investigation was not indicated by the evidence, but is within the doctrine, since the defendant's own state of mind and conduct satisfy the elements. An agent with neither evidence nor credence is outside both.

The third case deserves attention because it is where the doctrine reaches beyond the framework. An agent whose suspicion is unsupported by the available evidence, and who avoids confirming it, has not omitted an indicated investigation, so there is no violation $\nu_t>0$ in the framework's terms. The agent nevertheless has the motivational component that the framework treats as the aggravating feature of deliberate avoidance, and the law attributes knowledge to such an agent if the fact in question turns out to be true. Which of these treatments is more defensible is not settled by the formalism. It reflects a choice about whether culpability tracks the agent's own epistemic state and conduct, as the law's insistence on a subjective test suggests, or the agent's position relative to what the evidence supported, as the framework's account of duty suggests. The choice has consequences in the other direction as well, since the doctrine shields the agent who failed to notice evidence that was plainly there.

\section{Divergences}

The preceding section isolates one divergence, and several others can be stated more briefly. They are collected in Table~\ref{tab:diverge}, with the Canadian entries qualified by the limits of the sources.

\begin{table}[h]
\centering
\small
\renewcommand{\arraystretch}{1.25}
\begin{tabular}{p{2.1cm}p{3.7cm}p{3.5cm}p{3.5cm}}
\hline
Aspect & Framework & United States & Canada (secondary sources) \\
\hline
Trigger & Evidence supports the fact at or above $\theta_t$ & Subjective belief in a high probability & Suspicion rising to a perceived need for inquiry \\
Threshold & $\theta_t$, set by the standard of care, may depend on stakes and cost & $\eta$, fixed by the legal test & $\varsigma$, fixed by the legal test \\
Object & Materiality of what an investigation would reveal & Existence of a fact that is an element & Existence of a fact that is an element \\
Avoidance & Omission, coarser substitute, or interference & Deliberate actions to avoid learning & Deliberate decision not to inquire or not to inquire further \\
Effect & Classification of conduct as D1 or D2 & Substitutes for knowledge & Substitutes for knowledge \\
Competing constraints & Type J excludes culpability & Not part of the stated elements & Not part of the stated elements \\
Role and duty & Applicable set depends on role & Elements do not refer to role & Elements do not refer to role \\
Negligence & D1, a deficiency within the account & Expressly excluded & Expressly excluded \\
\hline
\end{tabular}
\caption{Points of comparison between the framework and the doctrine.}
\label{tab:diverge}
\end{table}

Four of these merit comment. First, the threshold. The framework's $\theta_t$ is fixed by the standard of care, and nothing in the account prevents the standard from requiring inquiry at a modest probability where the stakes are high and the investigation is cheap. The doctrine's level is a feature of the legal test and does not vary with the cost of inquiry or the gravity of what is at stake, at least as the sources state it. The consequence is that the framework can call for inquiry in circumstances where the law would find no willful blindness even if the agent looked away.

Second, avoidance. The American formulation speaks of deliberate actions to avoid learning, and the case on which it rests involved an interference. Whether a pure omission by a defendant who held a high-probability belief satisfies the second requirement is a question about the application of the formulation that the sources consulted do not resolve, and the paper does not claim that it is or is not satisfied. The Canadian formulation as reported speaks of a deliberate decision not to inquire, which on its face reaches omissions. The framework treats omission, coarser substitution, and interference alike as avoidance acts and distinguishes the motivated from the unmotivated by the motivational component and not by the form of the act.

Third, effect. The doctrine substitutes avoidance for knowledge, so that where the doctrine applies the defendant is treated as having known. The framework does not equate D2 with knowledge: an agent in D2 is deficient, culpably so, and the culpability is of avoidance and not of knowing. The difference is one of classification versus imputation, and Section~13 considers whether it matters.

Fourth, competing constraints. The framework recognizes that an omission can be justified by a norm that forbids the investigation, and it classifies such an omission as type J and not as deficient. The elements of the doctrine as stated do not refer to such a norm. How a court would treat an agent who declined to inquire because inquiry was forbidden, for instance by a confidentiality obligation or a statutory bar, is a question about whether the avoidance is deliberate in the relevant sense, and the sources consulted do not say.

\section{Recklessness, negligence, and willful blindness in two stages}

The law's trio of recklessness, negligence, and willful blindness can be recast using the two-stage structure of a decision, and doing so explains why the doctrine insists on keeping them apart. A decision has a first stage, the choice of investigation $q_t$, and a second stage, the choice of action $\pi_t$ after the observation. The justification relation $J_t$ evaluates both. A fault can lie in either stage.

Take the facts in the Court's description of the states of mind. A negligent defendant is one who should have known of a risk but in fact did not. In the model, the evidence indicated the investigation, the agent's credence did not reflect the risk, and the agent did not inquire; the fault is at the first stage, and it is a fault of omission without avoidance, which is type D1. A reckless defendant is one who knows of a substantial and unjustified risk. The agent's credence reflects the risk, no avoidance act is involved, and the fault lies in the second stage: the agent acts in the face of a risk that the agent perceives. A willfully blind defendant, on the doctrine, has a high credence and has avoided learning. The fault is at the first stage, in the choice of investigation or in what is allowed to flow into it, and it is accompanied by the credence and the motive.

\begin{definition}[Stylized states]
For agent $i$, fact $f$, and time $t$, let $\rho$ be a level of credence at which a risk is substantial. The agent is in the \emph{negligent state} if the investigation $q_f$ is indicated, $\beta^{(i)}_t(f)<\rho$, and no avoidance act is performed. The agent is in the \emph{reckless state} if $\beta^{(i)}_t(f)\ge\rho$, no avoidance act is performed, and the agent takes an action whose expected loss under the credence $\beta^{(i)}_t$ is substantially higher than that of an available alternative. The agent is in the \emph{willfully blind state} if $\beta^{(i)}_t(f)\ge\eta$, an avoidance act with respect to $f$ is performed, and $\mathrm{Mot}^{(i)}_t(f)$ holds.
\end{definition}

\begin{proposition}\label{prop:states}
Under these definitions the negligent and reckless states are mutually exclusive, and the willfully blind state is exclusive of both. The negligent state is a first-stage fault without avoidance, the reckless state is a second-stage fault without avoidance, and the willfully blind state is a first-stage fault with avoidance.
\end{proposition}

\begin{proof}
The negligent and reckless states are separated by the credence, which is below $\rho$ in the first and at or above it in the second. The willfully blind state requires an avoidance act, which the other two exclude by definition. The stage of fault is read from the definitions: the negligent and willfully blind states concern the choice of investigation, and the reckless state concerns the action chosen after the investigation.
\end{proof}

The proposition is a stylization and it is not a claim about how courts decide cases, since the three states of mind in the law are defined in prose and depend on context. What it provides is a reason for the doctrinal separation. The three states differ in where in the decision the fault lies and in whether an avoidance act is present, and the law's consequences attach differently to each. The doctrine's exclusion of negligence from willful blindness is therefore not a failure to see that the negligent agent should have inquired. It follows from the fact that the negligent agent's fault is a failure to acquire, with no avoidance, and that the law reserves the substitution of avoidance for knowledge to the agent who avoids.

\section{Partial inquiry, sham inquiry, and residual suspicion}

Both jurisdictions face the case in which the defendant did make some inquiry. The Canadian authority reported above holds that an inquiry does not by itself preclude willful blindness, and that the question is whether suspicion remained and further inquiry was avoided, whereas an agent who took all reasonable steps is not blind. The American case turns on an inquiry that was structured so as not to reveal the fact. The framework offers a way to state the difference between inquiries that discharge a duty and inquiries that do not.

Suppose agent $i$ performs an investigation $q$ with observation $Z_q$ and obtains outcome $z$, with posterior evidence-supported probability $p'=\PP(f\mid\F^{(i)}_t\vee\sigma(Z_q))$ evaluated at $z$. An investigation is a \emph{sham} with respect to $f$ if $Z_q$ is conditionally independent of $f$ given $\F^{(i)}_t$, and it is a \emph{restricted} investigation $q^{w}$ if it is the investigation $q$ performed with an input withheld, so that $\sigma(Z_{q^{w}})\subsetneq\sigma(Z_q)$. The withholding of the fact of copying from the patent attorney in the case discussed above is of this second kind: the attorney's study, performed on the full facts, would have been informative about whether the defendant's product infringed, and performed without the key fact it was less so.

\begin{proposition}\label{prop:inquiry}
(i) If $q$ is a sham with respect to $f$, then $p'=p^{(i)}_t(f)$ for every outcome, so the performance of $q$ leaves the indication of $q_f$ unchanged and does not discharge any duty to investigate $f$. (ii) The choice of a restricted investigation $q^{w}$ in place of $q$ is an avoidance act whenever $c_t(q^{w})\le c_t(q)$ and $q$ was available. (iii) If $p'\ge\theta_t$ and a refinement $q_2\in\Q_t^{\adm}$ is available, stopping is a deficient omission, of type D2 if $\mathrm{Mot}$ holds and D1 otherwise. (iv) If $p'<\theta_t$, the duty to investigate further lapses and stopping is not deficient.
\end{proposition}

\begin{proof}
(i) Conditional independence of $Z_q$ and $f$ given $\F^{(i)}_t$ gives $\PP(f\mid\F^{(i)}_t\vee\sigma(Z_q))=\PP(f\mid\F^{(i)}_t)$. (ii) The restricted investigation yields a coarser observation at no greater cost, which matches the definition of an avoidance act as the selection of a coarser substitute. (iii) and (iv) follow from the definition of indication applied at the updated information state and from the classification of deficient omissions.
\end{proof}

The Canadian rule, as reported, can be read in the same terms with one difference. It asks whether the accused \emph{remained suspicious}, and so it conditions on the credence after the inquiry, $\beta'$, and not on the evidence-supported probability $p'$. The framework conditions on $p'$. The two coincide for a Bayesian agent and may diverge for others for the reasons given in Section~6, and a person whose suspicion persists after a reassuring inquiry that the evidence in fact supports may be in a different position under the two accounts from a person whose suspicion lapses although the evidence does not warrant its lapse. The sham inquiry is the case in which the divergence is least troubling, since a defendant who arranges an inquiry that cannot reveal the fact has by that very arrangement shown what the credence was.

\section{Evidence and hindsight}

The framework's account of D2 includes a motivational component that, as the parent essay observed, cannot be inferred from the filtration or from the violation measure. The law faces the same problem in a more pressing form, since it must decide whether the component was present. The law's answer is circumstantial inference, and the circumstances on which the cases rely are instructive. The American case relies on copying that shows the defendant valued the design, on the defendant's knowledge of the industry practice that explains why the product bore no marking, and on the absence of any apparent reason for withholding the key fact from counsel other than to manufacture deniability. The Canadian authorities reported emphasize the preservation of deniability, the inherently suspicious character of the circumstances, and the contextual nature of the inquiry.

The structure of the inference can be stated in a way that explains why the American formulation requires deliberate actions. Let $H$ denote the hypothesis that the motivational component was present and $\neg H$ that the agent was negligent or innocent, and let the evidential force of an observation $E$ be the likelihood ratio $\Lambda(E)=\PP(E\mid H)/\PP(E\mid\neg H)$.

\begin{proposition}\label{prop:evidence}
(i) If a motivated agent and an unmotivated agent are equally likely to omit an inquiry, then the observation of an omission has $\Lambda=1$ and carries no evidence of motive. (ii) If a motivated agent is more likely than an unmotivated agent to perform an avoidance act of a given kind, then the observation of that act has $\Lambda>1$.
\end{proposition}

\begin{proof}
Both statements follow directly from the definition of $\Lambda$: equality of the two likelihoods gives a ratio of one, and a larger likelihood under $H$ gives a ratio exceeding one.
\end{proof}

The proposition is conditional, and its premises are empirical assumptions that the law makes implicitly. They are plausible for many omissions, since a negligent agent omits as readily as a motivated one, and for many affirmative acts of avoidance, since an unmotivated agent has no reason to arrange a consultation so that it cannot reveal what is at issue. The point of the proposition is that the evidential value of the stylized facts differs in the way the doctrine's emphasis suggests. A requirement of deliberate action to avoid learning is not merely a restriction on the scope of the doctrine. It also selects, among the facts that might be proved, those that discriminate between the motivated and the merely negligent, because a bare omission cannot do so.

Hindsight is a separate hazard. The doctrine's subjective test is anchored at the time of the act: it asks what the accused actually suspected and did then. A trier of fact who knows the outcome faces the temptation that the framework identifies as the confusion of the historical justification $J_t$ with the later assessment $\widehat J_{t\mid t+1}$. The outcome, and later statements and conduct, may properly be used as circumstantial evidence about what the accused knew and intended earlier, since they bear on the circumstances of the decision. They may not be used as though they were information the accused held when choosing. The doctrinal line between suspicion and knowledge coincides with the framework's line between what an agent's information supported at $t$ and what became known afterwards, and the discipline that the framework prescribes for later assessment is the discipline that a careful trier of fact must apply.

\section{Organizations}

The first companion paper argued that a pipeline can be locally justified at every stage while the organization loses information that was jointly indicated, that the organization is then answerable for whether mechanisms existed by which an indicated deficiency could reach someone with authority to act, and that an organization can be ignorant by arrangement without any individual looking away \cite{flyxion2026inputs}. Organizational avoidance, in its terms, requires an arrangement that suppresses adverse signals, knowledge on the part of an actor with authority over the arrangement that it has that effect, and persistence when an admissible mechanism that would remove the suppression was available. An organizational omission with the first but without the second and third is negligent and not avoidant.

The law addresses the attribution of mental states to organizations by rules of its own, and the rules differ between the two jurisdictions in ways that bear on this structure. The account here is limited to what the secondary sources consulted support.

In the United States, one line of authority permits the knowledge of several agents to be combined to establish the knowledge of a corporation, so that a company cannot escape liability by dividing duties among employees such that none of them knows enough. \emph{United States v.\ Bank of New England, N.A.}, 821 F.2d 844 (1st Cir.\ 1987), is reported in the commentary consulted to have involved a bank convicted of currency transaction reporting violations where it had failed to investigate a customer who withdrew sums above the reporting threshold by means of several checks and not one, the failure being treated as reflecting indifference amounting to willful blindness; the court, year, and reporter citation here are supplied from the author's recollection and are not confirmed by the sources consulted. The commentary adds that a corporation cannot avoid liability by declining to learn of its agents' illegal conduct.

In Canada, the Criminal Code provides in sections 22.1 and 22.2, as enacted in 2003 and as reproduced in the compilation consulted, separate rules for offences of negligence and for offences requiring a fault other than negligence. For a negligence offence, an organization is a party where a representative, or several representatives acting together, would be a party to the offence, and the senior officer responsible for the relevant part of the organization's activities, or the senior officers collectively, depart markedly from the standard of care that could reasonably be expected to prevent a representative from becoming a party. For an offence requiring a different fault, an organization is a party where a senior officer, acting at least in part to benefit the organization, commits the offence with the required mental state, directs other representatives so that they commit it while having that state, or knows that a representative is or is about to commit it and fails to take all reasonable measures to stop the representative. The definitions of representative and senior officer in the Code were not examined, and the in-force text should be confirmed.

The structures can be compared with the framework's. The American aggregation doctrine evaluates indication at the joint information $\F^{\vee}$ of the organization, in the notation of the companion paper, and attributes the result to the organization as knowledge. That is a stronger move than the framework's closure condition, which holds the organization answerable not for possessing the joint information but for maintaining mechanisms by which what is jointly indicated would reach an authorized actor. The framework would therefore support aggregation to the extent that an admissible mechanism would have achieved closure, and it would resist aggregation where no admissible mechanism could have brought the information together. This is a point about what the framework implies, and it is not a description of how the doctrine is applied. A doctrine that aggregates without that limit risks attributing to the organization knowledge that no one had and that no admissible arrangement could have produced.

The Canadian structure maps onto the framework's distinction between organizational negligence and organizational avoidance with some precision. Section 22.1, with its collective departure of senior officers from a standard of care and no requirement of a particular mental state, has the structure of an organizational omission of type D1. Section 22.2, which requires a senior officer's mental state or a senior officer's knowledge and failure to take reasonable measures, has the structure of the second and third conditions of organizational avoidance: knowledge at the level of an actor with authority, and persistence when reasonable measures were available. Read in this way, the proposition of the companion paper that an organization can be ignorant by arrangement without individual blindness corresponds to a situation in which the organization may fall within a negligence offence under the first rule and would fall outside an offence of subjective fault under the second unless a senior officer's state of mind can be shown. The law's division of attribution rules thus mirrors the framework's division of organizational omissions, and the fit suggests that the allocation is not arbitrary.

\section{A worked case}

A hypothetical shows how the layers separate. A distributor of electronic components receives a lot at a price far below the market from a new supplier. The fact at issue is $f$, that the lot is counterfeit, and a hypothetical provision attaches liability to dealing in counterfeit goods with knowledge that they are counterfeit. The actors are an account manager $M$, a compliance officer $C$, a purchasing clerk $P$, and a receiving clerk $R$. Five variants are considered, and in each the question put is what the framework, the American formulation, and the Canadian formulation would say on the stated facts, with the limits of the sources borne in mind.

In the first variant, $M$ regards it as probable that the lot is counterfeit, given the price and gaps in the paperwork, and tells $C$ only that the paperwork is delayed, omitting the price, so as to avoid a hold on the shipment. In the second, $M$ telephones the supplier, receives a reassurance that $M$ regards as self-serving, and stops verifying because further checking would delay shipment, while remaining doubtful. In the third, $M$ overlooks the price anomaly, which anyone in $M$'s role would have noticed, has no suspicion, and ships. In the fourth, $P$ sees the price and takes it to be a clearance sale, $R$ sees the paperwork gap and takes it to be clerical, there is no route by which either observation reaches the other or reaches $C$, and no one suspects. In the fifth, $M$ suspects counterfeiting, but the only way of verifying would require disclosing another customer's confidential information in breach of an enforceable undertaking, and $M$ ships without verifying.

\begin{table}[h]
\centering
\small
\renewcommand{\arraystretch}{1.25}
\begin{tabular}{p{1.2cm}p{3.6cm}p{3.6cm}p{3.6cm}}
\hline
Variant & Framework & American formulation & Canadian formulation \\
\hline
1 Interference & D2: indicated, admissible, interference, motive & Both elements met on the stated facts; the withholding resembles the paradigm case & Suspicion aroused and deliberate decision not to inquire properly, met on the stated facts \\
2 Partial inquiry & D2 if motive; stopping with $p'\ge\theta_t$ is deficient & Whether stopping short is a deliberate action to avoid learning is not resolved by the sources & Question whether suspicion remained and further inquiry was avoided; the reported rule contemplates it \\
3 Negligent & D1: indicated, no credence, no avoidance & Outside the doctrine; negligence expressly excluded & Outside the doctrine; not constructive knowledge \\
4 Fragmented & Individual $\nu=0$; organizational closure fails, D1 unless avoidance conditions hold & No individual meets the elements; aggregation doctrine is the route, subject to the limit proposed above & No individual meets the elements; s.\,22.1 for negligence offences, s.\,22.2 requires a senior officer's state of mind \\
5 Prohibited inquiry & J at the first stage; the second-stage action under a known risk may be reckless & Elements silent on the bar; effect on whether avoidance is deliberate unresolved & Elements silent on the bar; effect on whether avoidance is deliberate unresolved \\
\hline
\end{tabular}
\caption{Five variants of one hypothetical, classified under the framework and the two formulations.}
\label{tab:variants}
\end{table}

Table~\ref{tab:variants} records the comparisons, and several features deserve notice. The first and third variants show the doctrinal division of labor most clearly. The facts are identical but for the credence and the avoidance act, and the framework classifies the first as D2 and the third as D1, while the law treats the first as within willful blindness and the third as outside it, so that the line the law draws coincides with the framework's line between motivated and unmotivated deficiency. The second variant shows that the two accounts can diverge on the sequential case, since the framework asks whether the evidence-supported probability after the inquiry remained above threshold and the Canadian rule as reported asks whether the accused remained suspicious, and neither jurisdiction's formulation as stated in the sources consulted settles whether stopping is itself an avoidance act under the American text.

The fourth variant displays the organizational gap discussed in Section~11. Every individual is locally justified, so no individual-level element is satisfied, and the organization is answerable under the framework for whether a mechanism existed by which the two observations could have been combined. Whether the law reaches the organization depends on the attribution rule and on whether the offence is one of negligence or of subjective fault.

The fifth variant shows the value of the two-stage structure. At the first stage, the inquiry is forbidden by a competing norm, so the omission to inquire is of type J and is not deficient, and the framework does not treat $M$ as blind. At the second stage, $M$ acts in the face of a suspicion that the lot is counterfeit, which is a different fault and the one the law calls recklessness. The doctrine of willful blindness addresses the first stage, and the fault in the fifth variant, if there is one, lies in the second. A reading of the law that conflates the two would miss that an agent can be justified in not inquiring and unjustified in proceeding.

The retrospective investigation follows the same lines in each variant. The trier of fact must reconstruct what $M$ believed and did at the time, using what the record shows about $M$'s communications, the supplier's responses, and the routing of information in the firm. The later discovery that the lot was counterfeit is admissible as evidence about the circumstances, for instance that the anomaly was real, but it cannot supply $M$'s credence at the time. The variants differ in what the record can supply: in the first the unexplained omission to $C$ is the evidence of motive that Section~10 showed to be most probative, in the third no record shows any suspicion, and in the fourth the record may show that no actor held the combined information, which is the reason the organizational question is the live one.

\section{Should the law reach further}

The doctrine rests on an equivalence, expressed in the cases from \emph{Jewell} to \emph{Global-Tech}, between the person who knows and the person who has taken pains not to know. The Court in \emph{Global-Tech} reasoned that persons who know enough to blind themselves to direct proof of critical facts are in effect in the position of those with actual knowledge (at 766, quoting \emph{Jewell}). The dissent denied the equivalence and held that willful blindness is not knowledge and that the choice to treat it as such belongs to the legislature. The framework permits a more discriminating view of the equivalence than either side offers.

The framework's category D2 is defined by three conditions: the evidence indicated the fact, the investigation was admissible, and the agent omitted or interfered with it because of what the agent suspected it would show. An agent in D2 has the practical advantages of proceeding without the consequences of knowing, and the choice to avoid is an exercise of the will on behalf of the wrongful course. The equivalence with knowing wrongdoing is most plausible there, because what the agent has done is to arrange that the facts cannot be pressed on the agent while retaining the benefit of the course the facts would have forbidden. It is less plausible in the third combination of Section~6, in which the suspicion was not supported by the evidence, since the agent's avoidance then has no counterpart in the framework's account of what the agent owed. And it is not an argument for extending the doctrine to D1, since the negligent agent has neither the credence nor the avoidance on which the equivalence depends.

The framework therefore supports the law's exclusion of negligence from willful blindness, and the reason it supplies is that the two have different structures: negligence is a failure of acquisition without avoidance, willful blindness an act of avoidance, and the consequences that attach to knowledge are reserved for the latter. It also explains why criticism of the doctrine for failing to reach the ignorant organization or the careless individual tends to miss. Those cases are failures at the first layer of philosophical responsibility, to be addressed by duties of inquiry, standards of care, and rules of organizational attribution, and not by enlarging the second layer. Whether the existing duties and attribution rules are adequate is a separate question, which this paper does not take up. It is enough to say that the proper remedy for institutional ignorance lies in making the arrangement answerable, and not in treating the ignorant as knowing.

\section{Limits and conclusion}

Several limits should be stated plainly. The legal analysis is confined to federal law of the United States and to Canadian criminal law and is based on a limited set of sources. Of those, only the opinion in \emph{Global-Tech} was read in the original, and the Canadian decisions were consulted through secondary accounts that report paragraph numbers and characterizations without the text. The statements of those decisions are therefore at the level of generality that the accounts support, and any proposition resting on a precise formulation should be checked against the primary text before it is relied on. The statutory provisions were read in a compilation, and their in-force text and definitions were not confirmed. Other jurisdictions, including the states of the United States and the United Kingdom, have their own formulations and have not been considered.

The formal apparatus also has limits. The translation between a legal fact $f$ and the framework's materiality of an omitted variable is loose, since the law's elements concern the existence of a fact and the framework's indication concerns what an investigation would reveal. The legal levels $\eta$ and $\varsigma$, the diligence threshold $\theta_t$, and the substantiality level $\rho$ of the stylized states are all parameters that the account takes as given and does not derive. The stylized states of Section~8 are a model and are not a claim about how courts reason, and the evidential proposition of Section~10 depends on empirical premises about how motivated and unmotivated agents behave. The account says nothing about the standards of proof that govern a particular proceeding.

The conclusion is that the two uses of willful blindness are related but distinct, and that their relation can be stated exactly. The doctrine and the framework overlap where indicated evidence, a high credence, and deliberate avoidance coincide. They diverge because the doctrine turns on what the agent actually believed and did and the framework on what the evidence supported and the agent was bound to investigate, because the doctrine substitutes avoidance for knowledge where the framework classifies it, because the doctrine is silent on competing constraints and on role, and because the doctrine excludes negligence while the framework treats negligence as a deficiency within one account. The three layers of philosophical responsibility, legal knowledge, and evidential proof are separate functions with separate arguments, and the apparent disputes between them are largely disputes about which layer is in view. The division between possessed and acquirable information that organizes the philosophical account appears in the law as a division between doctrines, and the law's reasons for keeping them apart, which are that their consequences differ and that the avoidant are more culpable than the careless, are reasons the framework can accept.

\appendix
\section{Sources and the basis on which each was used}

Table~\ref{tab:sources} records, for each authority relied on, what it was used for and how its text was consulted. The distinction matters because several propositions about the Canadian law are stated only as reported.

\begin{table}[h]
\centering
\small
\renewcommand{\arraystretch}{1.25}
\begin{tabular}{p{4.3cm}p{4.3cm}p{4.6cm}}
\hline
Authority & Used for & Basis \\
\hline
\emph{Global-Tech Appliances, Inc.\ v.\ SEB S.A.}, 563 U.S.\ 754 (2011) & Two requirements; distinction from recklessness and negligence; facts; dissent & Opinion text read in the original (Library of Congress PDF) \\
\emph{United States v.\ Jewell}, 532 F.2d 697 (9th Cir.\ 1976) (en banc) & Origin of the American doctrine & As characterized in \emph{Global-Tech}; not read \\
Model Penal Code, definition of knowledge & High-probability formulation & As described in \emph{Global-Tech} (proposed draft); not read \\
\emph{United States v.\ Alston-Graves}, 435 F.3d 331 (D.C.\ Cir.\ 2006) & District of Columbia Circuit exception & As cited in \emph{Global-Tech}, n.\,9; not read \\
\emph{R.\ v.\ Sansregret}, [1985] 1 S.C.R.\ 570 & Contrast with recklessness & Secondary summary only \\
\emph{R.\ v.\ Jorgensen}, [1995] 4 S.C.R.\ 55 & Core test; contextual determination & Secondary summary only \\
\emph{R.\ v.\ Briscoe}, 2010 SCC 13 & Substitute for knowledge; limited scope; author and offences & Secondary summaries (annotated criminal law notebook; commentary reporting paragraph 23); primary text not retrieved \\
\emph{R.\ v.\ Lagace}, 2003 CanLII 30886 (Ont.\ C.A.); \emph{R.\ v.\ Farmer}, 2014 ONCA 823; \emph{R.\ v.\ Callejas}, 2011 ONCA 393; \emph{R.\ v.\ Malfara}, 2006 ONCA; \emph{R.\ v.\ Anderson}, 2020 ONCA 780 & Inquiry made; deliberate ignorance; not constructive knowledge; circumstantial proof & Secondary summary only \\
Criminal Code, R.S.C.\ 1985, c.\ C-46, ss.\ 22.1, 22.2 & Organizational liability & Text as reproduced in the UNODC SHERLOC database; in-force version and definitions not checked \\
\emph{United States v.\ Bank of New England, N.A.}, 821 F.2d 844 (1st Cir.\ 1987) & Collective knowledge; failure to investigate & Secondary summary only; court, year, and citation not confirmed by it \\
\hline
\end{tabular}
\caption{Authorities and the basis of use.}
\label{tab:sources}
\end{table}

\begin{thebibliography}{9}
\bibitem{flyxion2026domain} Flyxion. \emph{The Domain and Its Establishment: Information Acquisition and the Boundary of Responsibility}. October 2026.
\bibitem{flyxion2026inputs} Flyxion. \emph{Correct Relative to Its Inputs: Responsibility for the Domain of Automated Decisions}. October 2026.
\end{thebibliography}

\end{document}
